\documentclass{article}
\usepackage{amsmath}
\usepackage{amssymb}
\usepackage{amsfonts}
\usepackage[parfill]{parskip}
\title{Joint estimation of $\sigma_{m}^{2}$, $\lambda$ and $\sigma_{e}^{2}$ from population-level data}
\author{Sam Morris}
\date{October 2026}

\begin{document}


\maketitle

\section{Estimands}

We are interested in estimating three parameters from data on mutation counts and ages.

\begin{enumerate}
    \item $\sigma^2_m$ - variance per new mutation
    \item $K$ - the fitness landscape around a particular trait
    \item $\sigma^2_e$ - the residual variation
\end{enumerate}

\section{Model}

As before, we consider a standardised phenotype $y$ with variance $v_i$ given as:

\begin{equation}
    v_i = \sigma^2_m + \sigma^2_e + \sum_{k,m} C_{ikm} w(t_k, n_m)
\end{equation}

Where we have $k \times t$

We begin by assuming that the effect size of new mutations are drawn from a normal distribution $\beta \sim \mathcal{N}(0, \sigma^2_m)$ and that stabilising selection acts upon new mutations such that $s = K  \beta^2$. Our aim is to derive the weight function:

\begin{equation}
    w(t, n) = \frac{\mathbb{E}[\beta^2 | t, n]}{\sigma^2_m}
\end{equation}

Which tells us - compared to a brand new mutation, what is the expected effect size of a mutation with allele acount $n$ and age $t$.

We can rearrange this using Bayes' rule:

\begin{align}

\end{align}
